% ----- CHAUPAI 32 -----
\subsection*{\deva{द्वात्रिंशत्तम चौपाई} (Thirty-Second Chaupai)}
\addcontentsline{toc}{subsection}{\deva{द्वात्रिंशत्तम चौपाई} (Thirty-Second Chaupai) :  \deva{राम रसायन...}}

\begin{minipage}[t]{0.55\textwidth}
\vspace{0pt}

{\large \linespread{1.5}\selectfont
\deva{राम रसायन तुम्हरे पासा।}\\
\deva{सदा रहो रघुपति के दासा॥}
\par}

\vspace{0.5em}

{\large \linespread{1.5}\selectfont
\telu{రామ రసాయన తుమ్హరే పాసా।}\\
\telu{సదా రహో రఘుపతి కే దాసా॥}
\par}

\vspace{0.5em}

{\large \linespread{1.3}\selectfont
\textit{rāma rasāyana tumhare pāsā |}\\
\textit{sadā raho raghupati ke dāsā ||}
\par}
\end{minipage}%
\hfill
\begin{minipage}[t]{0.42\textwidth}
\vspace{0pt}
\begin{tcolorbox}[anchorbox]
\textbf{The Servant's Heart:} By consciously choosing to remain a humble servant and lifelong learner (\textit{Daasa}) rather than demanding to be the master, he found the ultimate cure (\textit{Rasayana}) for ego and unrest. Humility is the greatest medicine for a peaceful mind.
\vspace{0.5em}\newline He discovered that the ultimate cure for ego and unrest is choosing to remain a humble, lifelong learner.\newline\null\hfill\href{https://www.valmiki.iitk.ac.in/sloka?field_kanda_tid=6&language=dv&field_sarga_value=128}{\textit{Yuddha Kanda, 128:104}}
\end{tcolorbox}
\end{minipage}

\vspace{0.5em}
\subsubsection*{\textbf{Visual Rhythm Map:}}
\begin{table}[h!]
\centering
\renewcommand{\arraystretch}{1.2}
\setlength{\tabcolsep}{0pt}
\begin{tabularx}{\textwidth}{|*{16}{>{\centering\arraybackslash\hsize=1.0\hsize}X|}}
\hline
\multicolumn{3}{|>{\centering\arraybackslash\hsize=3.0\hsize}X|}{\textbf{\guru{रा}\laghu{म}} \par (3)} &
\multicolumn{5}{>{\centering\arraybackslash\hsize=5.0\hsize}X|}{\textbf{\laghu{र}\guru{सा}\laghu{य}\laghu{न}} \par (5)} &
\multicolumn{4}{>{\centering\arraybackslash\hsize=4.0\hsize}X|}{\textbf{\laghu{तु}\laghu{म्ह}\guru{रे}} \par (4)} &
\multicolumn{4}{>{\centering\arraybackslash\hsize=4.0\hsize}X|}{\textbf{\guru{पा}\guru{सा}} \par (4)} \\
\hline
\end{tabularx}
\vspace{-1.5em}
\end{table}

\begin{table}[h!]
\centering
\renewcommand{\arraystretch}{1.2}
\setlength{\tabcolsep}{0pt}
\begin{tabularx}{\textwidth}{|*{16}{>{\centering\arraybackslash\hsize=1.0\hsize}X|}}
\hline
\multicolumn{3}{|>{\centering\arraybackslash\hsize=3.0\hsize}X|}{\textbf{\laghu{स}\guru{दा}} \par (3)} &
\multicolumn{3}{>{\centering\arraybackslash\hsize=3.0\hsize}X|}{\textbf{\laghu{र}\guru{हो}} \par (3)} &
\multicolumn{4}{>{\centering\arraybackslash\hsize=4.0\hsize}X|}{\textbf{\laghu{र}\laghu{घु}\laghu{प}\laghu{ति}} \par (4)} &
\multicolumn{2}{>{\centering\arraybackslash\hsize=2.0\hsize}X|}{\textbf{\guru{के}} \par (2)} &
\multicolumn{4}{>{\centering\arraybackslash\hsize=4.0\hsize}X|}{\textbf{\guru{दा}\guru{सा}} \par (4)} \\
\hline
\end{tabularx}
\vspace{-1.5em}
\end{table}

\vspace{1em}
\textbf{ \deva{सरल-संस्कृतम्}:}\\
 \deva{भवतः पार्श्वे "राम-नाम" रूपी रसायनम् (औषधम्) अस्ति।}\\
 \deva{भवान् सर्वदा रघुपतेः (श्रीरामस्य) दासः भूत्वा एव तिष्ठति॥}\\


You possess the ultimate elixir, the essence of Lord Rama.\\
May you always remain the devoted servant of Lord Rama.


\vspace{1em}
\newpage
\textbf{\deva{प्रतिपदार्थः} (Meaning of each word):}
\begin{table}[h!]
\centering
\renewcommand{\arraystretch}{1.2}
\begin{tabularx}{\textwidth}{@{} l l X X @{}}
\toprule
\textbf{Awadhi} & \textbf{Sanskrit} & \textbf{English} & \textbf{Grammar \& Notes} \\
\midrule
\deva{राम रसायन} / \telu{రామ రసాయన}  &  \deva{राम-रसायनम्}  &  Elixir of Rama  & Rasa (Essence) + Ayana (Path) \\
\deva{तुम्हरे पासा} / \telu{తుమ్హరే పాసా}  &  \deva{भवतः पार्श्वे}  &  In your possession  &   \\
\deva{सदा रहो} / \telu{సదా రహో}  &  \deva{सर्वदा भव / तिष्ठ}  &  Always remain  &   \\
\deva{रघुपति के} / \telu{రఘుపతి కే}  &  \deva{रघुपतेः}  &  Of Lord Rama  &   \\
\deva{दासा} / \telu{దాసా}  &  \deva{दासः / सेवकः}  &  Servant  &   \\
\bottomrule
\end{tabularx}
\end{table}



\newpage
